% Regression: an atom threaded on a wire stands at no cell, so no endpoint
% names it and every face looks free.  Its carrier runs through the bead,
% and the two faces that run crosses are the ones the automatic station has
% to avoid.
\documentclass{standalone}
\usepackage{tenkz}
\makeatletter
\ExplSyntaxOn
\file_input:n { tenkz-kernel.code.tex }
\int_new:N \l__tenkz_test_side_count_int
\tl_new:N \l__tenkz_test_side_label_tl
\cs_new_eq:NN
  \__tenkz_test_side_auto:nN \__tenkz_kernel_r_label_auto_side:nN
\cs_set_protected:Npn \__tenkz_kernel_r_label_auto_side:nN #1#2
  {
    \__tenkz_test_side_auto:nN {#1} #2
    \int_gincr:N \l__tenkz_test_side_count_int
    \__tenkz_model_get:nnN {#1} {label} \l__tenkz_test_side_label_tl
    \str_case:VnF \l__tenkz_test_side_label_tl
      {
        {D}
          {
            \str_if_eq:VnF #2 {e}
              { \errmessage{a~bead~on~a~column~wire~kept~its~south} }
          }
        {E}
          {
            \str_if_eq:VnF #2 {s}
              { \errmessage{a~bead~on~a~row~wire~lost~its~free~south} }
          }
      }
      { \errmessage{unexpected~label~in~on-wire~fixture} }
  }
\ExplSyntaxOff
\makeatother
\begin{document}
\tenkzkernel
% A carrier down the column: north and south carry it, east is free.
\begin{tenkz}[lattice={3x1}, bonds=none]
  \tn[at=(1,1), name=P]{}
  \tn[at=(3,1), name=Q]{}
  \tnwire[name=W]{P}{Q}
  \tn[at=on W 0.5, skin=dot]{D}
\end{tenkz}
% A carrier along the row: east and west carry it, south stays free.
\begin{tenkz}[lattice={1x3}, bonds=none]
  \tn[at=(1,1), name=R]{}
  \tn[at=(1,3), name=S]{}
  \tnwire[name=Y]{R}{S}
  \tn[at=on Y 0.5, skin=dot]{E}
\end{tenkz}
\ExplSyntaxOn
\int_compare:nNnF { \l__tenkz_test_side_count_int } = {2}
  { \errmessage{on-wire~fixture~did~not~station~two~labels} }
\ExplSyntaxOff
\end{document}
